\documentclass[11pt]{article}
\usepackage[T1]{fontenc}
\usepackage[utf8]{inputenc}
\usepackage{amsmath,amssymb,amsthm}
\usepackage[margin=2.5cm]{geometry}
\usepackage{booktabs}
\usepackage{graphicx}
\setlength{\parindent}{1.4em}
\setlength{\parskip}{0pt}

\newtheorem{theorem}{Theorem}
\newtheorem{definition}{Definition}
\newtheorem{proposition}{Proposition}
\newtheorem{remark}{Remark}

\begin{document}

% ===== PAGE 1 =====
\noindent{\LARGE\textbf{Money or Vibe}}\\[0.3em]
\noindent{\small Why two rational people cannot agree on how to live.}\\[0.1em]
\noindent{\small Hakvin Vosteen \hfill 2026}
\vspace{0.3cm}

\begin{center}
\includegraphics[width=0.78\textwidth]{figures/hook_fig.pdf}
\end{center}

\vspace{1cm}

\medskip\noindent A banker thinks the artist is irrational. The artist thinks the banker is empty. Both are correct from inside their own utility function. Both are wrong about each other. The reason is structural, not about values or character.

\medskip\noindent The argument you are about to read is that humans optimize at least two distinct assets that have no exchange rate between them. Once you see this, the eternal money-versus-meaning conversation stops being a moral question and becomes a microeconomic one.

\newpage

% ===== PAGE 2 =====
\begin{definition}[Two-Asset Utility]
A person's utility depends on at least two distinct assets:
\[
U = U(M, V)
\]
where $M$ is Money Capital (transferable, accumulable, measurable) and $V$ is Vibe Capital (non-transferable, depreciable, context-specific).
\end{definition}

\noindent The key claim is that $V$ is not reducible to $M$. A rich person without Vibe Capital is genuinely poor in $V$. A Vibe-rich person without Money Capital is genuinely poor in $M$. They are not measuring the same thing.

\begin{theorem}[Inter-Personal Conflict from Incommensurable Utilities]
Consider person $A$ who maximizes $U_A = \alpha M + (1-\alpha) V$ with $\alpha \approx 1$, and person $B$ who maximizes $U_B = \beta M + (1-\beta) V$ with $\beta \approx 0$. Then there exists no exchange rate $\rho$ such that
\[
\Delta M = \rho \Delta V
\]
makes both indifferent. Therefore $A$ and $B$ cannot reach consensus on optimal allocation.
\end{theorem}

\begin{proof}
Suppose such a $\rho$ exists. Then for any $\Delta M$, $A$'s utility change is
\begin{align}
\Delta U_A &= \alpha \Delta M + (1-\alpha) \Delta V \\
&= \alpha \Delta M + (1-\alpha) \cdot \frac{\Delta M}{\rho}
\end{align}

For $A$ indifferent, $\Delta U_A = 0$:
\begin{align}
\alpha \Delta M + (1-\alpha) \cdot \frac{\Delta M}{\rho} &= 0 \\
\rho_A &= \frac{1-\alpha}{-\alpha}
\end{align}

By the same calculation for $B$:
\begin{align}
\rho_B &= \frac{1-\beta}{-\beta}
\end{align}

Since $\alpha \neq \beta$, we have $\rho_A \neq \rho_B$. The required exchange rate differs by person. There is no $\rho$ that satisfies both. $\hfill\blacksquare$
\end{proof}

\medskip\noindent The result is that Money-People and Vibe-People are not arguing about the same objective function. They live in different optimization spaces. Their disagreement is not solvable by sharing more information. It is structural.

\newpage

% ===== PAGE 3 =====
\begin{center}
\includegraphics[width=0.78\textwidth]{figures/utility_fig.pdf}
\end{center}

\medskip\noindent The figure shows two indifference curve families. The Money-Person treats high-$M$ low-$V$ positions as equivalent to mid-$M$ mid-$V$ positions. The Vibe-Person does the opposite. They literally see the same world differently.

\medskip\noindent\textbf{Why neither side can convert the other.}\quad The Money-Person says: with enough money, you can buy a good life, including relationships and meaning. Empirically false. Relationships need time, not money. Personal depth needs hours of reflection. Subcultural authenticity needs lived experience. None of these are buyable directly from $M$.

\medskip\noindent The Vibe-Person says: money does not matter if you live authentically. Empirically also problematic. Financial insecurity destroys Vibe through chronic stress. Material need forces performance mode. Without $M$-floor, $V$-accumulation collapses.

\begin{remark}[Mutual Floor Requirement]
Both extreme positions are unsustainable. Money-People need some $V > 0$ to avoid existential emptiness. Vibe-People need some $M > 0$ to avoid stress collapse. Both sides need minimum amounts of the other asset, but maximize in different directions. The disagreement is about what to do above the floor.
\end{remark}

\newpage

% ===== PAGE 4 =====
\begin{center}
\includegraphics[width=0.78\textwidth]{figures/conflict_fig.pdf}
\end{center}

\medskip\noindent\textbf{Why the conversation never resolves.}\quad The classic Arrow-Sen result that interpersonal utility comparison is impossible takes concrete form here. When a Money-Person says "you should focus on your career," they translate the recommendation through their own utility function. The Vibe-Person hears it through theirs. The same words mean different things.

\medskip\noindent\textbf{Concrete examples this explains.}
\begin{itemize}
\setlength\itemsep{0.2em}
\item Parents (money-frame) and art-studying kid (vibe-frame) argue endlessly about career choice
\item Investment banker and his Berlin artist friend never agree on how to live
\item Bay-Area tech-bro and NYC creative and Berlin bohemian argue in different languages
\item Influencer career vs investment-banking career are different asset-maximization strategies
\item Rich person who "is not happy" has M-maximum and V-deficit
\item Vibe-rich "broke but happy" person has the opposite
\end{itemize}

\medskip\noindent Each pair argues, neither convinces the other, and both walk away thinking the other is irrational. The framework predicts exactly this. They are not failing to communicate. They are succeeding at communicating, and discovering that they live in different optimization spaces.

\[
\boxed{\,\text{No } \rho \text{ exists with } \Delta M = \rho \Delta V \text{ that satisfies both} \;\Rightarrow\; \text{conflict is structural, not personal}\,}
\]

\newpage

% ===== PAGE 5 =====
\noindent\textbf{What follows from this}

\medskip\noindent The Money-versus-Vibe argument is not a moral disagreement that can be resolved by deeper conversation. It is a structural feature of human utility being multi-dimensional. Trying to win the argument is mistaking the problem.

\medskip\noindent What works instead: recognizing the asset-allocation logic both parties are running. Then the question becomes not "who is right" but "what mix is sustainable for me." Some people optimize $M$-heavy and live with their $V$-cost. Some optimize $V$-heavy and live with their $M$-cost. Some try to balance both, paying coordination costs but capturing both rewards.

\medskip\noindent None of these strategies is irrational. All have trade-offs. The only failure mode is being inside one strategy while not knowing the other exists.

\vspace{0.5cm}
\noindent\textbf{Notes on the framework}

\noindent The two-asset framework here extends Becker (1976), Akerlof and Kranton (2000), and Bourdieu (1979). The new claim is the Stock-Flow operationalization plus the incommensurability theorem. Vibe Capital depreciates without investment, accumulates with deliberate practice, and has no exchange rate to Money Capital.

\noindent Tradition for the underlying observation: Aristotle's megalopsychos, Nietzsche's Vornehmheit, Castiglione's sprezzatura, Veblen's instinct of workmanship, Bourdieu's habitus, Goffman's presentation of self. The phenomenon is two thousand years old. The economic formalization is new.

\noindent References available in working paper. Empirical strategies include cross-country surveys of utility-function dispersion, generational conflict-pattern analysis, and longitudinal lifecycle studies of asset accumulation.

\end{document}
